
\section{Entanglement, Helicity, and the Masslessness of the Photon}
\label{sec:entanglement_masslessness}

\subsection*{Overview}
One of the most subtle---and often misunderstood---arguments supporting a 
massless photon comes from the structure of quantum entanglement itself. 
Entangled photons exhibit correlations that match exactly those of a 
two-state quantum system. This turns out to be deeply connected to a 
fundamental property of massless spin-1 particles: they possess only two 
physical helicity states. Any non-zero photon mass would immediately alter 
these correlations in ways that have never been observed.

\subsection*{Entanglement Does Not Transmit Physical Influence}
When two photons are entangled, a measurement on photon~A does not induce 
any physical change in photon~B. No signal, force, or energy travels between 
them. Instead, the two photons share a single non-factorizable quantum 
state. Measurement collapses our \emph{information} about the system, not 
the physical state of the distant photon. Thus, entanglement does not imply 
superluminal influence and does not require masslessness in itself.

\subsection*{Helicity Structure of a Massless Spin-1 Particle}
The crucial point is different: a massless spin-1 particle such as the 
photon has only two physical degrees of freedom,
\[
\lambda = +1, \qquad \lambda = -1,
\]
corresponding to the two transverse helicity states.  
This reduction from three to two degrees of freedom is only possible if
\[
m_\gamma = 0.
\]
A massive spin-1 particle always carries three polarization states,
\[
\lambda = +1, \qquad \lambda = 0, \qquad \lambda = -1,
\]
where the additional $\lambda=0$ state represents the longitudinal mode.
This longitudinal component cannot be removed without imposing 
gauge symmetry, and gauge symmetry is incompatible with a Proca mass term.

\subsection*{Why Entanglement Tests Reveal the Spin Structure}
Experiments with entangled photons---including Bell tests, polarization 
correlation measurements, Hong--Ou--Mandel interference, and 
quantum-state tomography---all rely sensitively on the number of physically 
accessible polarization states. The observed correlations match perfectly 
those predicted by a two-state (qubit-like) structure. 

If the photon had even an exceedingly small mass, the longitudinal mode 
would contribute to the entangled state. This would modify:
\begin{itemize}
	\item the angular dependence of Bell-type correlations,
	\item the visibility of polarization interference,
	\item the structure of two-photon coincidence rates,
	\item the transformation properties under Lorentz boosts, and
	\item the reconstruction of the density matrix in tomography.
\end{itemize}
None of these deviations have ever been observed.

\subsection*{Absence of a Longitudinal Mode as Evidence for Masslessness}
The empirical fact that all entanglement experiments with photons are 
consistent with exactly two polarization states constitutes a strong 
indirect constraint on the photon mass. A longitudinal degree of freedom 
would be experimentally unmistakable. Its absence supports:
\[
m_\gamma \approx 0
\quad\text{with extremely tight bounds from laboratory and astrophysical data}.
\]

\subsection*{Conclusion}
Entanglement itself does not transmit physical influence between distant 
photons. Instead, it reveals which polarization states the photon 
\emph{can} possess. The experimentally verified two-state helicity 
structure is only compatible with a massless spin-1 field. Any non-zero 
photon mass would introduce a third polarization mode and fundamentally 
alter the entanglement signatures we observe.

Thus, the remarkably consistent two-state behavior of entangled photons 
provides compelling indirect evidence that the photon is, to extremely 
high precision, massless.
